\begin{abstract}
Recent phase-ordering methods summarise pairwise pass interactions in a single graph learned from
a training corpus and apply it to every program. We test that assumption in LLVM with a
pre-registered study. We measure \NPairs{} ordered pass-pair interactions on \NPrograms{} programs
at \NSnapshots{} points along LLVM's \texttt{-O2} pipeline (\NExperiments{} experiments, none
failed). Three of four pre-registered hypotheses are supported.
(1) Interactions are program-specific. Ordering with a program's own graph instead of a static
one reduces IR size by \HOnePct\% at a budget of ten compilations. The gain carries over to
\texttt{.text} size on x86-64 (\MSizeXPct\%) and AArch64 (\MSizeAPct\%).
(2) Interactions drift. The median rank correlation between interactions on raw IR and after
inlining is \HTwoRho, and the largest change happens in early simplification.
(3) Interactions are predictable. A model trained on other programs removes \HThreeGap\% of the
static graph's error.
(4) Re-predicting interactions along the pipeline does not beat a genetic algorithm at ten
compilations. Exploratory analysis shows that graph-based sequence builders plateau even with the
program's own measured graph, while a predicted graph beats every search method when only one
compilation is allowed. One graph does not fit all programs. The next step is to use per-program
predictions inside search rather than as a fixed graph.
\end{abstract}
